\chapter*{Registro del documento: lo que no se pudo verificar}
\addcontentsline{toc}{chapter}{Registro: lo que no se pudo verificar}

Este registro es parte del documento y no una nota técnica. Un capítulo que dice con precisión
qué no puede afirmar es mejor que uno que lo afirma sin fuente.

\section*{Uno. Lo que es imposible verificar bajo régimen de lectura en vivo}

Las tres son de la misma familia: exigen el presupuesto \emph{aprobado} del mismo ejercicio,
que el día de la entrega del paquete no existe.

\begin{list}{}{\setlength{\leftmargin}{1.4em}\setlength{\itemsep}{3pt}}
\item[1.] \textbf{Separar objeto de perímetro} en cualquier discrepancia contra otro
documento. Cuando dos lecturas del mismo paquete difieren, la diferencia puede ser de método
---qué se sumó--- o de momento ---proyecto contra aprobado---, y sin el aprobado no se
distinguen.
\item[2.] \textbf{Medir lo que cambió la Cámara} entre el proyecto y el presupuesto que se
ejerce.
\item[3.] \textbf{Adjudicar un titular del proyecto} contra el presupuesto efectivamente
aprobado. En el ejercicio anterior, la Cámara aumentó el gasto educativo lo suficiente para
cambiar el signo del capítulo entre noviembre y diciembre.
\end{list}

\textbf{No son carencias de esta corrida: son la definición del régimen.} Todo documento
publicado a las setenta y dos horas de la entrega tiene estas tres limitaciones, las declare o
no.

\section*{Dos. Denominadores que no se obtuvieron}

Se buscaron en los datos abiertos de las instituciones responsables y en el catálogo nacional,
dos veces cada uno.

\begin{list}{}{\setlength{\leftmargin}{1.4em}\setlength{\itemsep}{2pt}}
\item[$\bullet$] \textbf{Matrícula por nivel educativo.} Sin ella no se publica gasto por
alumno y no se separa el efecto demográfico del de cobertura (capítulo 6).
\item[$\bullet$] \textbf{Población derechohabiente de IMSS y de ISSSTE.} Sin ella no se
publica gasto por afiliado (capítulo 5).
\item[$\bullet$] \textbf{Padrón de pensionados por régimen.} Sin él no se publica pensión
media (capítulo 7).
\end{list}

\textbf{Lo que sí se publicó} son los denominadores que sí se obtuvieron ---población total,
población de 65 años y más, población en edad escolar, padrón de personas sin seguridad social
de IMSS-Bienestar--- en cuadros separados de los presupuestales y con su fuente.

\section*{Tres. Piezas del paquete que el servidor no entrega}

\textbf{La exposición de motivos del proyecto de presupuesto lleva cinco ejercicios caída.} Las
páginas índice del portal la enumeran y los archivos devuelven error. Es el documento donde el
Ejecutivo explica sus decisiones de gasto, y su ausencia afecta a este documento en tres
lugares concretos: el perímetro canónico de salud, la explicación de la aportación de capital a
Pemex y las estrategias programáticas por ramo.

\section*{Cuatro. Cifras que el paquete publica y no permite descomponer}

\begin{list}{}{\setlength{\leftmargin}{1.4em}\setlength{\itemsep}{2pt}}
\item[$\bullet$] \textbf{La meta tributaria}, entre actividad económica, eficiencia
recaudatoria y las tres leyes de la miscelánea.
\item[$\bullet$] \textbf{La aportación de capital a Pemex}, entre amortización de pasivos e
inversión productiva.
\item[$\bullet$] \textbf{La compensación por inflación} sobre la deuda: el paquete no
identifica qué índice de precios le corresponde y el efecto neto de la inflación sobre el
acervo queda en un rango de $-0,05$ a $-2,10$ puntos del PIB.
\item[$\bullet$] \textbf{La fracción de cada programa} que corresponde a cada anexo
transversal.
\item[$\bullet$] \textbf{La tasa del derecho petrolero}: la iniciativa de 2026 no contiene el
artículo que la fijaba entre 2020 y 2023. Se buscó y no está.
\end{list}

\section*{Cinco. Un residuo que este documento no nombra}

La reconciliación entre los impuestos del artículo 1o.\ de la iniciativa de ingresos
---5\,838\,541,1 millones--- y los ingresos tributarios del cuadro de finanzas públicas
---5\,838\,571,0--- deja \textbf{29,9 millones de pesos} sin explicar. Es cinco diezmilésimas
de por ciento y se registra porque una diferencia sin nombre es una diferencia sin nombre.

\section*{Seis. Lo que este documento hizo mal y corrigió}

Se registra porque el propósito de este ejercicio es aprender a leer un paquete en vivo, y los
errores propios enseñan más que los aciertos.

\begin{list}{}{\setlength{\leftmargin}{1.4em}\setlength{\itemsep}{3pt}}
\item[1.] \textbf{Se filtró por nombre y devolvió cero.} El primer cálculo del gasto de las dos
empresas del Estado dio cero para las dos, porque cambiaron de denominación en los analíticos.
Es la misma trampa que esta casa lleva años señalando y volvió a caer en ella.
\item[2.] \textbf{Se filtró por clave completa y perdió un programa.} La Pensión Mujeres
Bienestar es U316 en un ejercicio y S316 en el siguiente: cambió la modalidad, no el número. La
regla «por clave, nunca por nombre» resultó insuficiente y quedó corregida a «por el número de
programa, verificando la modalidad contra el diff».
\item[3.] \textbf{Se teclearon cifras de memoria, dos veces.} Los renglones de la iniciativa de
ingresos del ejercicio anterior y el monto de un ramo. En los dos casos las cifras eran falsas y
se detectaron al cotejar contra la fuente, antes de que llegaran al texto. \textbf{Toda cifra
entra al archivo de datos desde su fuente o no entra.}
\end{list}

\section*{Siete. Cómo se validó lo que sí se afirma}

Treinta y tres identidades contables y de coherencia entre anexos, todas cerradas, corridas
después de cada cambio en los datos. El gasto bruto del proyecto ---11\,746\,796,8 millones de
pesos--- se obtuvo por tres rutas independientes en dos servidores distintos y coincide al peso
con el Anexo 1 del decreto. El total de la iniciativa de ingresos coincide al peso con el gasto
neto total. La identidad entre el flujo y el acervo cierra con un residuo de 5\,999,7 millones,
quince milésimas de punto del PIB.

\textbf{Se valida por identidad y no por relectura.} Un dígito mal leído rompe una suma; una
relectura lo repite.
\clearpage
